\songtitle{Güiro quemado}

\tag{original-source}

\begin{notes}
A smooth jazz-tinged song that emphasizes phrasing, atmosphere, and lyrical intimacy. %
The lyrics begin with “Y es de noche”, giving a quick sense of the song’s tone.
\end{notes}

\begin{style}
salsa, afro-cuban jazz, 180 BPM, male lead vocal, smooth clear vocals, piano tumbao, trombone stabs, brass shout chorus, güiro syncopation, tambora accents, handclap montuno, call and response chorus, minor key montuno, dynamic breakdowns, live room reverb, tape saturation, mono low end, bittersweet longing, urgent release
\end{style}

\begin{lyrics}

\songpart{Intro}
\annotation{piano enters with a mysterious tumbao, the Caribbean güiro and the brass instruments with a sharp hit}

\annotation{spoken, piano only}\\
Y es de noche\\
una noche larga e interminable\\
donde la vigilia es un sueño\\
en donde los recuerdos se agolpan\\
atropellando a su paso\\
la poca autoestima que aún se tiene.

\songpart{Verse 1}
Y es de noche todavía, larga noche interminable,\\
con una vigilia horrible que sobre mi ser porfía.\\
La memoria en agonía va atropellando el sentido,\\
con recuerdo malherido, la poca fe que aún retengo,\\
y mientras despierto vengo voy quedándome perdido.

\songpart{Chorus}
Y es noche sobre mi vida, y el recuerdo vuelve a entrar,\\
con la memoria vencida que no me deja escapar.

\annotation{Instrumental Break, güiro and brass}

\songpart{Verse 2}
Hay pensamientos extraños que no deben razonarse,\\
ni detenerse a mirarse como si evitaran daños.\\
Pasan dejando sus años sobre la vida y la mente,\\
capturando lentamente una razón sin razón,\\
y admiro sólo la ilusión de no pensar de repente.

\songpart{Chorus}
Callo lo que necesito, mientras el momento huye,\\
y en silencio me marchito viendo el instante que fluye.

\annotation{Instrumental Break, tambora and cymbals}\\
\annotation{shouted}
¡Salsa! ¡Que suene el piano! ¡Que rujan los trombones que el sujeto calla!\\
\annotation{piano and brass interlude}

\songpart{Verse 3}
Es una espera infinita, eterna y desesperada,\\
espera desamparada que nunca se debilita.\\
Y aunque la noche se agita la sigo buscando aún,\\
sin rastro, señal ni alud, extraño palabras viejas,\\
una imagen que ya alejas desde tu invisible luz.

\songpart{Chorus}
Y es noche sobre mi vida, y el recuerdo vuelve a entrar,\\
con la memoria vencida que no me deja escapar.

\songpart{Montuno}
\annotation{Call and Response, energetic, call and response style}\\
\annotation{Choir:} ¿Qué es lo que te preocupa?\\
\annotation{Solo:} Que es noche sobre mi vida.\\
\annotation{Choir:} ¿Qué te preocupa de eso?\\
\annotation{Solo:} Que el recuerdo vuelve a entrar.

\annotation{Choir:} ¿Qué pasa cuando te callas?\\
\annotation{Solo:} En silencio me marchito.\\
\annotation{Choir:} ¿Qué es lo que tanto callas?\\
\annotation{Solo:} Callo lo que necesito.

\annotation{Choir:} ¿Cuándo es que te marchitas?\\
\annotation{Solo:} Mientras el momento huye.\\
\annotation{Choir:} ¿Cuándo es que lo necesitas?\\
\annotation{Solo:} Viendo el instante que fluye.

\songpart{Verse 4}
Es mi vida la que espera, la que depende al hablar,\\
y vuelvo siempre a callar cuando el momento se fuera.\\
Mi alma entera quisiera decirte cuánto te extraño,\\
pero el silencio hace daño y las palabras no llegan.\\
Y mis pensamientos juegan a escabullirse otro año.

\songpart{Chorus}
Y es noche sobre mi vida, y el recuerdo vuelve a entrar,\\
con la memoria vencida que no me deja escapar.

Callo lo que necesito, mientras el momento huye,\\
y en silencio me marchito viendo el instante que fluye.

\songpart{Outro}
\annotation{melodic, free verse, fading}\\
Un espacio. %
Una oportunidad.\\
Es de nuevo la ocasión esperada.\\
... %
y espero.

y calladamente vi escabullirse\\
esta última vez.

\end{lyrics}

